\subsection{Example: the case of dimension 1}

In this section, one considers an integral, noetherian ring A of dimension 1, admitting a dualizing module \(\Omega\) . We denote by K the field of fractions of A, and by V the K-vector space \(K \otimes_{A} \Omega\) .

The canonical homomorphism \(\Omega\to V\) is injective, and the K-vector space \(V\) is of dimension 1 ( \(n^{\circ}\) 2, cor. 1 of prop. 4); let us identify \(\Omega\) to a sub-A-module of \(V\) .

PROPOSITION 8. The A-module V/Ω is a Matlis module.

Consider the exact sequence

\[
0 \to \Omega \to \mathrm{V} \to \mathrm{V} / \Omega \to 0;
\]

the A-module V is injective (A, X, p.~18, example 1), and we have \(\mathrm{di}_{\mathrm{A}}(\Omega) = 1\) ( \(n^{\circ} 1\) , prop. 1). From this we deduce on the one hand that \(V/\Omega\) is injective ( \(\S 3\) , \(n^{\circ} 1\) , prop. 1), on the other hand, that for every maximal ideal \(\mathfrak{m}\) of A, the \(\Lambda/\mathfrak{m}\) -vector space \(\mathrm{Hom}_{\mathrm{A}}(\mathrm{A}/\mathfrak{m}, \mathrm{V}/\Omega)\) is isomorphic to \(\mathrm{Ext}_{\mathrm{A}}^{1}(\mathrm{A}/\mathfrak{m}, \Omega)\) , hence of dimension 1. Since \(V/\Omega\) is a torsion module, its associated prime ideals are maximal; this proves the proposition ( \(\S 8\) , \(n^{\circ} 4\) ).

Let M be an A-module; in accordance with loc. cit., we denote by D(M) the A-module \(\mathrm{Hom}_{\mathrm{A}}(\mathrm{M},\mathrm{V}/\Omega)\) . One can apply the constructions of no. 5 by taking for I the complex

\[
\dots 0 \longrightarrow \mathrm{V} \xrightarrow {p} \mathrm{V} / \Omega \longrightarrow 0 \dots
\]

where V is placed in degree 0, and p denotes the canonical surjection. The complex D(M) is

\[
\dots 0 \longrightarrow \operatorname{Hom} _ {\mathrm{A}} (\mathrm{M}, \mathrm{V}) \xrightarrow {p _ {\mathrm{M}}} \mathrm{D} (\mathrm{M}) \longrightarrow 0 \dots ,
\]

with \(p_{\mathbf{M}} = \mathrm{Hom}(\mathbf{1}_{\mathbf{M}}, p)\) We have a canonical isomorphism of graded A-modules \(\varphi(\mathbf{M}, \mathbf{I}) : \mathrm{H}(\mathbf{D}(\mathbf{M})) \longrightarrow \mathrm{Ext}_{\mathbf{A}}(\mathbf{M}, \Omega)\) (A, X, p.~100, th. 1).

When M is a torsion module, the module \(\mathbf{D}(\mathbf{M})^{0} = \operatorname{Hom}_{\mathbf{A}}(\mathbf{M}, \mathbf{V})\) is zero, and \(\varphi(\mathbf{M}, \mathbf{I})\) is an isomorphism from D(M) to \(\operatorname{Ext}_{\Lambda}^{1}(\mathbf{M}, \Omega)\) ; the morphism \(\alpha_{M}: M \to \mathbf{D}(\mathbf{D}(\mathbf{M}))\) is none other than the canonical homomorphism of A-modules

\[
\alpha_ {\mathrm{M}}: \mathrm{M} \longrightarrow \mathrm{D} (\mathrm{D(M)})
\]

defined in § 8, no. 3, which is an isomorphism when M is of finite type (that is, of finite length). In this case we recover the situation of loc. cit.

Let us return to the general case, and suppose the A-module M is of finite type. Then D(M) is a torsion module, so the A-module \(\mathbf{D}(\mathbf{D}(\mathbf{M}))^{-1} = \mathrm{Hom}_{\mathrm{A}}(\mathrm{D}(\mathrm{M}), \mathrm{V})\) is zero. On the other hand, the A-module \(\mathrm{Hom}_{\mathrm{A}}(\mathbf{D}(\mathbf{M})^{0}, \mathrm{V}) = \mathrm{Hom}_{\mathrm{A}}(\mathrm{Hom}_{\mathrm{A}}(\mathrm{M}, \mathrm{V}), \mathrm{V})\) is naturally identified with \(K \otimes_{A} M\) , so that the homomorphism

\[
\operatorname{Hom} (1, p): \operatorname{Hom} _ {\mathrm{A}} (\operatorname{Hom} _ {\mathrm{A}} (\mathrm{M}, \mathrm{V}), \mathrm{V}) \longrightarrow \mathrm{D} (\operatorname{Hom} _ {\mathrm{A}} (\mathrm{M}, \mathrm{V}))
\]

is identified with the map

\[
j: \mathrm{K} \otimes_ {\mathrm{A}} \mathrm{M} \longrightarrow \mathrm{D} (\operatorname{Hom} _ {\mathrm{A}} (\mathrm{M}, \mathrm{V}))
\]

such that \(j(\lambda \otimes m)(f) = p(\lambda f(m))\) for \(\lambda \in \mathbf{K}\) , \(m \in \mathbf{M}\) , \(f \in \mathrm{Hom}_{\mathrm{A}}(\mathrm{M},\mathrm{V})\) . Th. 1 of no. 5 therefore translates into the exactness of the sequence

\[
0 \longrightarrow \mathrm{M} \xrightarrow {(i , \alpha_ {\mathrm{M}})} (\mathrm{K} \otimes_ {\mathrm{A}} \mathrm{M}) \oplus \mathrm{D} (\mathrm{D} (\mathrm{M})) \xrightarrow {(j , - \mathrm{D} (p _ {\mathrm{M}}))} \mathrm{D} (\operatorname{Hom} _ {\Lambda} (\mathrm{M}, \mathrm{V})) \longrightarrow 0,
\]

where \(i\) denotes the canonical map from M into \(\mathrm{K} \otimes_{\mathrm{A}} \mathrm{M}\) The kernel of \(i\) is identified with the torsion submodule \(\mathrm{T(M)}\) of M, and its cokernel at \((\mathrm{K/A}) \otimes_{\mathrm{A}} \mathrm{M}\) . Consider the commutative diagram with exact rows

\[
\begin{array}{c c c c c c c} 0 \to \mathrm{T(M)} & \longrightarrow & \mathrm{M} & \stackrel {{\imath}} {{\longrightarrow}} & \mathrm {K\otimes_ {A} M} & \longrightarrow & (\mathrm{K/A}) \otimes_ {\mathrm{A}} \mathrm{M} \longrightarrow 0 \\ & & ^ {\alpha_ {\mathrm{M}}} \Bigg \downarrow & & \Bigg \downarrow_ {j} \\ 0 \to \mathrm{D(Ext} _ {\mathrm{A}} ^ {1} (\mathrm{M}, \Omega)) & \longrightarrow & \mathrm{D(D(M))} & \stackrel {{\mathrm{D} (p _ {\mathrm{M}})}} {{\longrightarrow}} & \mathrm{D(Hom} _ {\mathrm{A}} (\mathrm{M}, \mathrm{V})) & \longrightarrow & \mathrm{D(Hom} _ {\mathrm{A}} (\mathrm{M}, \Omega)) \longrightarrow 0 \end{array}
\]

where the second line is obtained by Matlis duality from the exact sequence

\[
0 \rightarrow \operatorname{Hom} _ {\mathrm{A}} (\mathrm{M}, \Omega) \longrightarrow \operatorname{Hom} _ {\mathrm{A}} (\mathrm{M}, \mathrm{V}) \xrightarrow {p _ {\mathrm{M}}} \operatorname{Hom} _ {\mathrm{A}} (\mathrm{M}, \mathrm{V} / \Omega) \longrightarrow \operatorname{Ext} _ {\Lambda} ^ {1} (\mathrm{M}, \Omega) \rightarrow 0.
\]

Theorem 1 then means that the homomorphisms of A-modules

\[
\gamma^ {0} (\mathrm{M}): \mathrm{T} (\mathrm{M}) \longrightarrow \mathrm{D} \left(\operatorname{Ext} _ {\mathrm{A}} ^ {1} (\mathrm{M}, \boldsymbol {\Omega})\right) e t \quad \gamma^ {1} (\mathrm{M}): (\mathrm{K} / \mathrm{A}) \otimes_ {\Lambda} \mathrm{M} \longrightarrow \mathrm{D} \left(\operatorname{Hom} _ {\mathrm{A}} (\mathrm{M}, \boldsymbol {\Omega})\right)
\]

inferred from \(\alpha_{\mathrm{M}}\) and \(j\) respectively, are bijective. Since the A-module T(M) has finite length, the A-module \(\operatorname{Ext}_{\mathrm{A}}^{1}(\mathrm{M},\Omega)\) is of finite length and is identified with the Matlis dual D(T(M)), and we have \([\operatorname{Ext}_{\mathrm{A}}^{1}(\mathrm{M},\Omega)] = [\operatorname{T}(\mathrm{M})]\) in the group Z0(A) and \(\operatorname{long}_{\mathrm{A}}(\operatorname{Ext}_{\mathrm{A}}^{1}(\mathrm{M},\Omega)) = \operatorname{long}_{\mathrm{A}}(\operatorname{T}(\mathrm{M}))\) ( \(\S 8\) (no. 4, prop. 4). On the other hand, when one takes M = A, one obtains a canonical isomorphism \(\gamma^{1}(\mathrm{A}):K / A\to D(\Omega)\) .

Let B be a subring of K containing A, finite over \(\Lambda\) . For every maximal ideal \(\mathfrak{m}\) of A, we have \(\mathrm{prof}_{\mathrm{A}_{\mathfrak{m}}}(B_{\mathfrak{m}}) = \dim_{A_{\mathfrak{m}}}(B_{\mathfrak{m}}) = 1\) ( \(\S 1\) , \(n^{\circ} 1\) , Remark 2 and VIII, \(\S 2\) , \(n^{\circ} 3\) , Th. 1), so that B is a Macaulay A-module. Consequently, the B-module \(\Omega_{\mathrm{B}} = \mathrm{Hom}_{\mathrm{A}}(B, \Omega)\) is dualizing \(n^{\circ} 3\) , remark). The canonical map from \(\Omega_{\mathrm{B}} = \mathrm{Hom}_{\mathrm{A}}(B, \Omega)\) in \(\Omega = \mathrm{Hom}_{\mathrm{A}}(A, \Omega)\) is injective; its image consists of the elements \(\omega\) of \(\Omega\) such that the sub-B-module B \(\omega\) of V be contained in \(\Omega\) Thus \(\Omega_{\mathrm{B}}\) is identified with the largest sub-B-module of \(\Omega\) . The A-module B/ \(\Lambda\) is of finite length; the exact sequence

\[
0 \to \Omega_ {\mathrm{B}} \to \Omega \to \operatorname{Ext} _ {\mathrm{A}} ^ {1} (\mathrm{B/A}, \Omega) \to 0
\]

allows one to identify \(\Omega/\Omega_{\mathrm{B}}\) with \(\operatorname{Ext}_{\mathrm{A}}^{1}(\mathrm{B}/\mathrm{A},\Omega)\) , hence by what precedes with \(\mathrm{D}(\mathrm{B}/\mathrm{A})\) . In particular, we have \([\mathrm{B}/\mathrm{A}] = [\Omega/\Omega_{\mathrm{B}}]\) in \(Z_{0}(\mathrm{A})\) and \(\operatorname{Ann}_{\mathrm{A}}(\mathrm{B}/\mathrm{A}) = \operatorname{Ann}_{\mathrm{A}}(\Omega/\Omega_{\mathrm{B}})\) ( \(\S 8\) , \(\mathfrak{n}^{\circ}4\) , prop. 4).

The ideal \(\mathfrak{c} = \mathrm{Ann}_{\mathrm{A}}(\mathrm{B / A})\) is the transporter A : B, that is (VII, § 1, no. 1) the set of elements \(x\) of K such that \(xB \subset \Lambda\) . It is a (nonzero) ideal of

A and of B; in fact it is the largest ideal of B contained in A. Since \(\Omega_{\mathrm{B}} : \Omega \subset \Omega : \Omega = \mathrm{A} (\mathrm{n}^{\circ} 4, \text{prop. } 7, \mathrm{b}))\) , we have Ann \(_{\mathrm{A}}(\Omega/\Omega_{\mathrm{B}}) = \Omega_{\mathrm{B}} : \Omega\) , whence finally

\[
\mathfrak {c} = \operatorname{Ann} _ {\mathrm{A}} (\mathrm{B/A}) = \operatorname{Ann} _ {\mathrm{A}} (\Omega / \Omega_ {\mathrm{B}}) = \Omega_ {\mathrm{B}}: \Omega .
\]

Since \(\Omega_{\mathrm{B}}\) is a B-module, the relation \(x\Omega \subset \Omega_{\mathrm{B}}\) is equivalent to \(xB\Omega \subset \Omega_{\mathrm{B}}\) , so that we also have \(\mathfrak{c} = \Omega_{\mathrm{B}}: \mathrm{B}\Omega\) .

We shall now specialize the preceding discussion to the case where B is the integral closure of A; the hypothesis that B is a finitely generated A-module is satisfied when the ring A is Japanese (IX, § 4, no. 1, def. 1), which is the case when it is local and complete (loc. cit., no. 2, th. 2), or when it is essentially of finite type over a field (loc. cit., no. 1, remark 2 and example). The ring B is then a Dedekind ring (VII, § 2, no. 2, th. 1), and the torsion-free B-modules \(\Omega_{\mathrm{B}}\) , B \(\Omega\) and c are projective of rank 1 (VII, § 4, no. 10, prop. 22). The relation \(c = \Omega_{\mathrm{B}} : \mathrm{B}\Omega\) then means that the linear map \(c \otimes_{\mathrm{B}} \mathrm{B}\Omega \to \Omega_{\mathrm{B}}\) deduced from the action of K on V is an isomorphism (II, § 5, no. 6, prop. 11). In particular, we have \(\Omega_{\mathrm{B}} = c(\mathrm{B}\Omega) = c\Omega\) .

PROPOSITION 9. Let B be the integral closure of A, and \(\mathfrak{c} = \mathrm{A}:\mathrm{B}\) . Suppose that B is a finitely generated A-module. We have the inequality \([\mathrm{B} / \mathrm{c}] \leqslant 2[\mathrm{B} / \mathrm{A}]\) in \(\mathrm{Z}_0(\mathrm{A})\) . For equality to hold, it is necessary and sufficient that A be a Gorenstein ring.

We have \([B/c] = [B/A] + [A/c]\) , so that the inequality under consideration is equivalent to \([A/c] \leqslant [B/A]\) .

\begin{enumerate}
\def\labelenumi{\Alph{enumi})}
\item
  For every maximal ideal m of A, the integral closure of \(A_{m}\) is \(B_{m}\) (V, § 1, no. 5, cor. 1), and we have \(c_{m} = A_{m} : B_{m}\) . Moreover, by definition we have \([B/c] = \sum_{m} \text{long}_{A_{m}}(B_{m}/c_{m})[m]\) and \([B/A] = \sum_{m} \text{long}_{A_{m}}(B_{m}/A_{m})[m]\) . This reduces us to proving the proposition when the ring A is local, which we shall henceforth assume. In this case the ordered group \(Z_{0}(A)\) is canonically identified with Z, so that the class of a module of finite length is its length. The ring B is semi-local and the B-module BΩ is free of rank 1 (II, § 5, no. 3, prop. 5).
\item
  If A is a Gorenstein ring, the A-module \(\Omega\) is free of rank 1 (no. 4, cor. 1 of prop. 7); let us choose a generator \(\omega\) of \(\Omega\) . One has \(\Omega = A\omega\) and \(\Omega_{\mathrm{B}} = c\Omega = c\omega\) and consequently \(\text{long}(A/c) = \text{long}(\Omega/\Omega_{\mathrm{B}}) = \text{long}(B/A)\) .
\item
  Suppose the residue field \(\kappa_{\mathrm{A}}\) is infinite. For every maximal ideal \(\mathfrak{n}\) of B, denote by L(n) the B/n-vector space (of dimension 1) BΩ/nBΩ, and by pr \(_{\mathfrak{n}}\) the canonical projection of \(\bigoplus_{\mathfrak{n}}\) L(n) onto L(n). Let \(\varphi : \Omega \longrightarrow \bigoplus_{\mathfrak{n}}\) L(n) be the restriction to \(\Omega\) of the canonical homomorphism BΩ \(\longrightarrow \bigoplus_{\mathfrak{n}}\) L(n). The image of \(\varphi\) is a sub- \(\kappa_{\mathrm{A}}\) -vector space of \(\bigoplus_{\mathfrak{n}}\) L(n); it is not contained in Ker pr \(_{\mathfrak{n}}\) , for otherwise one would have \(\Omega \subset \mathfrak{n}\mathrm{B}\Omega\) and consequently BΩ \(\subset \mathfrak{n}\mathrm{B}\Omega\) , which is contradictory. Thus the image of \(\varphi\) is not contained in the union of the Ker pr \(_{\mathfrak{n}}\) (A, V, p.~40, lemma 1); there therefore exists an element \(\omega\) of \(\Omega\) whose image in BΩ/nBΩ is nonzero for every n, which implies that \(\omega\) generates the B-module BΩ (II, § 3, no. 3, prop. 11).
\end{enumerate}

Let \(a \in \mathbf{A}\) ; if \(a\omega\) belongs to \(\Omega_{\mathrm{B}}\) , one has \(a\mathrm{B}\omega \subset \Omega_{\mathrm{B}}\) , hence \(a\Omega \subset \Omega_{\mathrm{B}}\) , which implies \(a \in \mathfrak{c}\) . The map \(a \mapsto a\omega\) therefore induces an injection of \(\mathrm{A}/\mathfrak{c}\) in \(\Omega/\Omega_{\mathrm{B}}\) ; consequently one has \(\mathrm{long}(\mathrm{A}/\mathfrak{c}) \leqslant \mathrm{long}(\Omega/\Omega_{\mathrm{B}}) = \mathrm{long}(\mathrm{B}/\mathrm{A})\) .

If long(A/c) = long(B/A), then Aω + ΩB = Ω. We may assume that the ideal c is contained in mA (otherwise A equals B, hence is a Gorenstein ring). Since ΩB = cΩ is contained in mAΩ, it follows from Nakayama's lemma that ω generates Ω. Thus the A-module Ω is cyclic, hence free of rank 1, which means that A is a Gorenstein ring (no. 4, cor. 1 of prop. 7).

\begin{enumerate}
\def\labelenumi{\Alph{enumi})}
\setcounter{enumi}{3}
\tightlist
\item
  Let us treat the general case. Denote by A' the ring A{]}X{[}, that is (IX, App., no. 2) the local ring of the polynomial ring Λ{[}X{]} in the prime ideal m\_AΛ{[}X{]} ; it is a flat, integral A-algebra of dimension 1, whose residue field κ\_A' is identified with κ\_A(X) and the field of fractions to K(X) (loc. cit.). By the corollary of Prop. 5 of No. 3, the A'-module A' ⊗\_A Ω is dualizing. Set B' = A' ⊗\_A B; this is the integral closure of Λ' in K(X) (V, § 1, no. 3, prop. 13 and no. 5, prop. 16). The transporter c' = A' : B' is equal to cA' (I, § 2, no. 10, formula (11)). For every A-module M of finite length, one has long\_A'(A' ⊗\_A M) = long\_A(M): indeed, since the A-algebra A' is flat, it suffices to prove this relation when M is simple, that is, isomorphic to κ\_A; but in this case A' ⊗\_A κ\_A is identified with κ\_A′, whence our assertion. We therefore have
\end{enumerate}

\[
\operatorname{long} _ {\mathrm{A}} (\mathrm{B} / \mathfrak {c}) = \operatorname{long} _ {\mathrm{A} ^ {\prime}} \left(\mathrm{B} ^ {\prime} / \mathfrak {c} ^ {\prime}\right) \quad \text { and } \quad \operatorname{long} _ {\mathrm{A}} (\mathrm{B} / \mathrm{A}) = \operatorname{long} _ {\mathrm{A} ^ {\prime}} \left(\mathrm{B} ^ {\prime} / \mathrm{A} ^ {\prime}\right).
\]

The ring A' satisfies the hypotheses of the proposition, and its residue field is infinite. By part C) of the proof, we have \(\text{long}_{\text{A}'}(\text{B}'/\text{c}') \leqslant 2 \text{long}_{\text{A}'}(\text{B}'/\text{A}')\) , and the equality implies that A' is a Gorenstein ring; but this latter condition entails that A is a Gorenstein ring (§ 3, no. 8, cor. 1 of prop. 12).

